% Options for packages loaded elsewhere
\PassOptionsToPackage{unicode}{hyperref}
\PassOptionsToPackage{hyphens}{url}
\PassOptionsToPackage{dvipsnames,svgnames,x11names}{xcolor}
\documentclass[
  english,
  11pt,
  reqno]{amsart}
\usepackage{xcolor}
\usepackage{amsmath,amssymb}
\setcounter{secnumdepth}{-\maxdimen} % remove section numbering
\usepackage{iftex}
\ifPDFTeX
  \usepackage[T1]{fontenc}
  \usepackage[utf8]{inputenc}
  \usepackage{textcomp} % provide euro and other symbols
\else % if luatex or xetex
  \usepackage{unicode-math} % this also loads fontspec
  \defaultfontfeatures{Scale=MatchLowercase}
  \defaultfontfeatures[\rmfamily]{Ligatures=TeX,Scale=1}
\fi
\usepackage{lmodern}
\ifPDFTeX\else
  % xetex/luatex font selection
\fi
% Use upquote if available, for straight quotes in verbatim environments
\IfFileExists{upquote.sty}{\usepackage{upquote}}{}
\IfFileExists{microtype.sty}{% use microtype if available
  \usepackage[]{microtype}
  \UseMicrotypeSet[protrusion]{basicmath} % disable protrusion for tt fonts
}{}
\makeatletter
\@ifundefined{KOMAClassName}{% if non-KOMA class
  \IfFileExists{parskip.sty}{%
    \usepackage{parskip}
  }{% else
    \setlength{\parindent}{0pt}
    \setlength{\parskip}{6pt plus 2pt minus 1pt}}
}{% if KOMA class
  \KOMAoptions{parskip=half}}
\makeatother
\usepackage{longtable,booktabs,array}
\newcounter{none} % for unnumbered tables
\usepackage{calc} % for calculating minipage widths
% Correct order of tables after \paragraph or \subparagraph
\usepackage{etoolbox}
\makeatletter
\patchcmd\longtable{\par}{\if@noskipsec\mbox{}\fi\par}{}{}
\makeatother
% Allow footnotes in longtable head/foot
\IfFileExists{footnotehyper.sty}{\usepackage{footnotehyper}}{\usepackage{footnote}}
\makesavenoteenv{longtable}
\ifLuaTeX
\usepackage[bidi=basic,shorthands=off]{babel}
\else
\usepackage[bidi=default,shorthands=off]{babel}
\fi
\ifLuaTeX
  \usepackage{selnolig} % disable illegal ligatures
\fi
\setlength{\emergencystretch}{3em} % prevent overfull lines
\providecommand{\tightlist}{%
  \setlength{\itemsep}{0pt}\setlength{\parskip}{0pt}}
% Portable additions for the arXiv-oriented profile.
% Keep packages within the standard TeX Live distribution.
\usepackage{mathtools}
\usepackage{amssymb}
\usepackage{amsthm}
\usepackage{booktabs}
\usepackage{listings}
\usepackage{microtype}
\usepackage{centernot}
\usepackage{newunicodechar}

% Portable negated hooked arrow used by several manuscripts.
\providecommand{\nhookrightarrow}{\centernot\hookrightarrow}

% U+220E END OF PROOF: render a small hollow square at the right margin.
\newcommand{\rightqedsquare}{%
  \ifhmode
    \unskip\nobreak\hfill\ensuremath{\square}%
  \else
    \noindent\hfill\ensuremath{\square}%
  \fi
}
\newunicodechar{∎}{\rightqedsquare}

% Reduce avoidable overfull lines without changing the source manuscript.
\setlength{\emergencystretch}{2em}
% Common LaTeX macros for CCG papers
% Keep this minimal and portable across engines.

% Math shortcuts
\newcommand{\N}{\mathbb{N}}
\newcommand{\Z}{\mathbb{Z}}
\newcommand{\Q}{\mathbb{Q}}
\newcommand{\R}{\mathbb{R}}
\newcommand{\C}{\mathbb{C}}
\newcommand{\eps}{\varepsilon}

% Theorem-like environments (lightweight)
% If you need full theoremstyles, consider amsthm or thmtools here.
% \usepackage{amsthm}   % already available via TeX Live if needed
% \newtheorem{theorem}{Theorem}
% \newtheorem{lemma}{Lemma}

% Code listings default (avoid minted to keep shell-escape off)
\lstset{basicstyle=\ttfamily\small,breaklines=true,columns=fullflexible}

\usepackage{bookmark}
\IfFileExists{xurl.sty}{\usepackage{xurl}}{} % add URL line breaks if available
\urlstyle{same}
\hypersetup{
  pdftitle={No Set Carries Exactly Two Dense Linear Orders without Endpoints},
  pdfauthor={Lior Isthmus},
  pdflang={en},
  colorlinks=true,
  linkcolor={blue},
  filecolor={Maroon},
  citecolor={blue},
  urlcolor={blue},
  pdfcreator={LaTeX via pandoc}}

\title{No Set Carries Exactly Two Dense Linear Orders without Endpoints}
\usepackage{etoolbox}
\makeatletter
\providecommand{\subtitle}[1]{% add subtitle to \maketitle
  \apptocmd{\@title}{\par {\large #1 \par}}{}{}
}
\makeatother
\subtitle{A Cut-Rotation Proof in a Weak Zermelo Theory without Choice
or Replacement}
\author{Lior Isthmus}
\date{}

\begin{document}
\maketitle

\markboth{\MakeUppercase{Lior Isthmus}}{\MakeUppercase{No Set Carries Exactly Two Dense Linear Orders without\ldots{}}}

\protect\phantomsection\label{sec:title}

\protect\phantomsection\label{sec:abstract}

\section*{Abstract}\label{abstract}

Let \(Z_{\mathrm{sep}}\) consist of Extensionality, Pairing, Infinity,
Union, Power Set, and the full Separation schema, and write \(s(X)=n\)
when \(X\) carries exactly \(n\) isomorphism types of dense linear
orders without endpoints. No form of Choice, Replacement, or Foundation
is assumed. We prove

\hypertarget{eq:abstract-main}{}

\begin{equation*}
Z_{\mathrm{sep}}
\vdash
\neg\exists X\,(s(X)=2).
\end{equation*}

Using countable-carrier uniqueness and the Dedekind-infinite four-type
alternative from the exact-three companion, exact two forces \(X\) to be
neither at most countable nor Dedekind-infinite and every DLO on \(X\)
to be rigid. A type-count-free dyadic reflection argument eliminates the
possibility that both types are self-dual, leaving one rigid
non-self-dual dual pair \(\{[L],[L^*]\}\).

For each \(c\in L\), four one-point cut orders form an antipodal
two-colored square. Vertical monochromaticity would self-dualize both
rays and then the whole carrier; hence

\[
L\cong\operatorname{Rot}_c(L).
\]

Rigidity gives a unique isomorphism
\(\rho_c:L\to\operatorname{Rot}_c(L)\). Define
\(\delta(c)=\rho_c^{-1}(c)\). Its graph is obtained by Separation inside
\(X\times X\), without forming a Replacement-generated family
\((\rho_c)_{c\in X}\). Opposite-ray exchange and hereditary
Dedekind-finiteness force \(\delta\) to be an injective strictly
decreasing surjection. Thus \(\delta\) is an order reversal of \(L\),
contradicting \(L\not\cong L^*\).

Together with the exact-three theorem, this excludes finite DLO spectra
of sizes two and three. Spectra of size at least four are not decided
here.

\textbf{Keywords.} dense linear orders; axiom of choice; Dedekind-finite
sets; cut rotation; categoricity; weak set theory.

\textbf{2020 Mathematics Subject Classification.} Primary 03E25;
Secondary 03E30, 03C35, 03C64, 06A05.

\tableofcontents

\protect\phantomsection\label{sec:introduction}

\section{1. Introduction}\label{introduction}

A \textbf{dense linear order without endpoints}, abbreviated
\textbf{DLO}, is a nonempty strict linear order \((X,<)\) which is dense
and has neither a least nor a greatest point. Throughout this paper,

\[
X\text{ is at most countable}
\quad\Longleftrightarrow\quad
X\hookrightarrow\omega,
\]

and

\[
X\text{ is Dedekind-infinite}
\quad\Longleftrightarrow\quad
\omega\hookrightarrow X.
\]

Without Choice these conditions need not be complementary.

In the language \(\{<\}\), the DLO axioms form the complete theory
\(\operatorname{Th}(\mathbb Q,<)\); see \hyperref[ref:marker2002]{Marker
2002}. Thus the finite spectra considered here are precisely the numbers
of models of this theory, up to internal isomorphism, on one fixed
underlying set. For background on cardinal comparison and
Dedekind-finite sets without Choice, see \hyperref[ref:jech1973]{Jech
1973} and \hyperref[ref:howardrubin1998]{Howard--Rubin 1998}.

The companion paper \hyperref[ref:isthmus2026a]{Isthmus 2026a} proves in
the same weak theory that no set carries exactly three DLO types,
thereby answering Shelah's Question 7.11 negatively
\hyperref[ref:shelah2009]{Shelah 2009}. Two global spectrum results are
used directly below: at-most-countable DLO uniqueness and the four-type
alternative on every Dedekind-infinite carrier. Several elementary
order-theoretic and foundational lemmas are also reused, including
rigidity facts, finite ranking, binary-word coding, finite block sums,
and bounded recursion. \hyperref[app:dyadic-engine]{Appendix B} of the
present paper extracts and proves the type-count-free dyadic reflection
engine needed for exact two, rather than appealing to an exact-three
theorem with altered hypotheses. The earlier Lean 4 development
formalizes the exact-three theorem package
\hyperref[ref:isthmus2026b]{Isthmus 2026b}.

A companion Lean 4 development formalizes all 45 numbered items in this
paper and the claimed proof-theoretic endpoint
\hyperref[ref:isthmus2026c]{Isthmus 2026c}. It includes the
object-language theory \(Z_{\mathrm{sep}}\), a natural-deduction
calculus, and the completeness bridge used to produce the final
derivation. Its final declaration is
\texttt{ExactTwoDLO.FO.zsep\_proves\_not\_spec2Sentence}; its type is
the derivation assertion

\[
\operatorname{Provable}\,Z_{\mathrm{sep}}
\bigl(\mathord\sim\operatorname{spec2Sentence}\bigr).
\]

The exact-three release is imported there as a pinned, read-only
dependency; the exact-two reflection engine and cut-rotation argument
are implemented in the separate \texttt{ExactTwoDLO} development.

Exact two has a different obstruction from exact three. On three types,
reversal has a fixed point and therefore supplies a self-dual type. On
two types, reversal may exchange the two types and have no fixed point.
The main task is therefore to break one rigid non-self-dual dual pair
without first assuming a self-dual representative.

The proof is organized around two statements.

\protect\phantomsection\label{thm:intro-reduction}

\textbf{Theorem 1.1 (Exact-two structural reduction).} If \(s(X)=2\),
then:

\begin{enumerate}
\tightlist
\item
  \(X\nhookrightarrow\omega\);
\item
  \(\omega\nhookrightarrow X\);
\item
  every DLO on \(X\) is rigid;
\item
  reversal exchanges the two types;
\item
  no DLO on \(X\) is self-dual;
\item
  after naming one representative \(L\), the two types are
\end{enumerate}

\[
   \{[L],[L^*]\},
   \qquad
   L\not\cong L^*.
\]

\protect\phantomsection\label{thm:intro-cut-reversal}

\textbf{Theorem 1.2 (Cut-preimage reversal).} Let \(L=(X,<)\) satisfy
the residual conditions of Theorem 1.1. If every one-point cut rotation
\(\operatorname{Rot}_c(L)\) is isomorphic to \(L\), then the unique
cut-rotation isomorphisms determine a decreasing bijection

\[
\delta:X\longrightarrow X.
\]

Consequently \(L\cong L^*\).

The finite cut square proves the hypothesis of Theorem 1.2. The
contradiction with Theorem 1.1 then excludes exact two.

The new mechanism is local. Exact-three exclusion uses a self-dual
representative and a full dyadic reflection tree. Exact-two exclusion
uses the tree only to eliminate the all-self-dual reversal pattern. Once
the remaining dual pair is isolated, a four-corner cut square and the
point-indexed preimage map \(\delta\) close the proof.

The paper is a genuine companion rather than a duplicate. The
foundational coding, countable benchmark theorem, and Lindenbaum theorem
are imported from \hyperref[ref:isthmus2026a]{Isthmus 2026a}. The
present paper proves the exact-two reduction, a parameterized
type-count-free reflection-tree engine, the cut-square theorem, and the
cut-preimage reversal.

\protect\phantomsection\label{sec:formal-setting}

\section{2. Formal setting and imported
results}\label{formal-setting-and-imported-results}

\protect\phantomsection\label{def:2-1-base-theory}

\textbf{Definition 2.1 (The base theory).} \(Z_{\mathrm{sep}}\) consists
of Extensionality, Pairing, Infinity, Union, Power Set, and the full
Separation schema. No use is made of Choice, Countable Choice, Dependent
Choice, Replacement, Collection, or Foundation.

All orders are subsets of \(X\times X\), and all witnessing maps are
internal sets. The construction of \(\omega\), bounded unique recursion,
finite functions, finite block sums, and the relevant arithmetic are
available in this theory by \hyperref[ref:isthmus2026a]{Isthmus 2026a,
Appendices A--B}.

\protect\phantomsection\label{def:2-2-spectrum}

\textbf{Definition 2.2 (Finite spectrum formulas).} For every standard
finite \(n\ge1\), let \(\operatorname{Spec}_{\ge n}(X)\) assert that
there are \(n\) pairwise nonisomorphic DLO relations on \(X\). Put

\begin{equation}\protect\phantomsection\label{eq:finite-spectrum-formula}{
\operatorname{Spec}_{=n}(X)
:=
\operatorname{Spec}_{\ge n}(X)
\wedge
\neg\operatorname{Spec}_{\ge n+1}(X).
}\end{equation}

The notation \(s(X)\ge n\) abbreviates
\(\operatorname{Spec}_{\ge n}(X)\), and \(s(X)=n\) abbreviates
\(\operatorname{Spec}_{=n}(X)\). No quotient set of isomorphism classes
is formed.

If \(s(X)=n\) and \(R_0,\ldots,R_{n-1}\) are pairwise nonisomorphic DLOs
on \(X\), then every DLO on \(X\) is isomorphic to exactly one \(R_i\).
Otherwise an additional type would witness \(s(X)\ge n+1\).

\protect\phantomsection\label{def:2-3-order-notation}

\textbf{Definition 2.3 (Order notation).} For sets \(A,B\), the notation
\(A\hookrightarrow B\) means that an injection from \(A\) to \(B\)
exists, and \(A\nhookrightarrow B\) means that no such injection exists.
For \(m\in\omega\), the notation \(|A|\le m\) abbreviates the existence
of an injection \(A\to m\), and \(|A|=m\) abbreviates the existence of a
bijection \(A\to m\).

For a linear order \(L=(X,<)\):

\begin{itemize}
\tightlist
\item
  \(L^*\) is the reverse order;
\item
  \((-\infty,c)_L\) and \((c,\infty)_L\) are the strict rays at \(c\);
\item
  an \textbf{order reversal} of \(L\) is an anti-isomorphism \(L\to L\);
\item
  \(L\) is \textbf{self-dual} if \(L\cong L^*\);
\item
  \(L\) is \textbf{rigid} if its increasing automorphism group is
  trivial.
\end{itemize}

Finite order sums are taken over disjoint blocks. Empty blocks are
omitted, and every new adjacent boundary is checked against the finite
block-sum DLO criterion of \hyperref[ref:isthmus2026a]{Isthmus 2026a,
Lemma 2.7}.

\protect\phantomsection\label{thm:2-4-subcountable-uniqueness}

\textbf{Theorem 2.4 (At-most-countable DLO uniqueness).} If \(X\)
carries a DLO and \(X\hookrightarrow\omega\), then

\begin{equation}\protect\phantomsection\label{eq:subcountable-uniqueness}{
s(X)=1.
}\end{equation}

\emph{Proof.} This is \hyperref[ref:isthmus2026a]{Isthmus 2026a, Lemma
3.11}. The proof first shows that an infinite subcountable set is
bijective with \(\omega\), then carries out a deterministic least-index
back-and-forth construction inside a fixed power set. ∎

\protect\phantomsection\label{thm:2-5-dedekind-alternative}

\textbf{Theorem 2.5 (Dedekind-infinite four-type alternative).} If \(X\)
carries a DLO and \(\omega\hookrightarrow X\), then

\begin{equation}\protect\phantomsection\label{eq:dedekind-alternative}{
X\hookrightarrow\omega
\quad\text{or}\quad
s(X)\ge4.
}\end{equation}

\emph{Proof.} This is \hyperref[ref:isthmus2026a]{Isthmus 2026a, Theorem
3.12}. A choice-free monotone-subsequence construction gives a countable
benchmark whose DLO complement is either subcountable or supports four
pairwise nonisomorphic same-carrier DLOs. ∎

\protect\phantomsection\label{lem:2-6-automorphism-orbit}

\textbf{Lemma 2.6 (A nontrivial increasing automorphism gives a
countable orbit).} If a linear order on \(X\) has a nonidentity
increasing automorphism, then

\begin{equation}\protect\phantomsection\label{eq:automorphism-orbit}{
\omega\hookrightarrow X.
}\end{equation}

\emph{Proof.} This is \hyperref[ref:isthmus2026a]{Isthmus 2026a, Lemma
6.1}. Choose \(x<g(x)\) after replacing \(g\) by \(g^{-1}\) if
necessary, and recursively form the strictly increasing orbit
\(g^n(x)\). ∎

\protect\phantomsection\label{cor:2-7-common-host}

\textbf{Corollary 2.7 (Common-host self-duality criterion).} If one
linear order contains initial segments isomorphic to both \(A\) and
\(A^*\), or final segments isomorphic to both, then

\begin{equation}\protect\phantomsection\label{eq:common-host}{
A\cong A^*.
}\end{equation}

\emph{Proof.} This is \hyperref[ref:isthmus2026a]{Isthmus 2026a,
Corollary 4.3}, obtained from Lindenbaum's initial/final-segment
theorem; see also \hyperref[ref:ervin2017]{Ervin 2017, §2.3}. ∎

\protect\phantomsection\label{lem:2-8-rigidity-convex}

\textbf{Lemma 2.8 (Rigidity and reversals on convex suborders).} Let
\(L\) be rigid.

\begin{enumerate}
\tightlist
\item
  Every convex suborder of \(L\) is rigid.
\item
  Every reversal of a rigid order is involutive.
\item
  A rigid self-dual order has a unique reversal.
\end{enumerate}

\emph{Proof.} This is \hyperref[ref:isthmus2026a]{Isthmus 2026a, Lemma
4.8}. A convex-suborder automorphism extends by the identity, two
reversals differ by an automorphism, and the square of a reversal is an
automorphism. ∎

\protect\phantomsection\label{sec:structural-reduction}

\section{3. Exact-two structural
reduction}\label{exact-two-structural-reduction}

Throughout this section assume

\[
s(X)=2,
\]

and fix pairwise nonisomorphic representatives

\[
L_0=(X,<_{0}),
\qquad
L_1=(X,<_{1}).
\]

\protect\phantomsection\label{lem:3-1-not-subcountable}

\textbf{Lemma 3.1 (An exact-two carrier is not at most countable).} One
has

\[
X\nhookrightarrow\omega.
\]

\emph{Proof.} Otherwise Theorem 2.4 would give \(s(X)=1\). ∎

\protect\phantomsection\label{lem:3-2-not-dedekind-infinite}

\textbf{Lemma 3.2 (An exact-two carrier is not Dedekind-infinite).} One
has

\[
\omega\nhookrightarrow X.
\]

\emph{Proof.} If \(\omega\hookrightarrow X\), Theorem 2.5 gives either
\(X\hookrightarrow\omega\) or \(s(X)\ge4\). The first alternative
contradicts Lemma 3.1. In the second, discarding one of four witnesses
gives \(s(X)\ge3\), contradicting \(s(X)=2\). ∎

\protect\phantomsection\label{lem:3-3-rigidity}

\textbf{Lemma 3.3 (Every exact-two DLO is rigid).} Every DLO on \(X\)
has trivial increasing automorphism group.

\emph{Proof.} A nontrivial increasing automorphism would give
\(\omega\hookrightarrow X\) by Lemma 2.6, contradicting Lemma 3.2. ∎

\protect\phantomsection\label{lem:3-4-hereditary-dedekind}

\textbf{Lemma 3.4 (Hereditary Dedekind-finiteness).} Let
\(Y\subseteq X\). Every injective map

\[
j:Y\longrightarrow Y
\]

is surjective.

\emph{Proof.} Suppose \(j\) were injective and not surjective. Choose

\[
y_0\in Y\setminus j[Y]
\]

and define

\[
y_{n+1}=j(y_n)
\]

by bounded unique recursion inside the fixed set \(\omega\times Y\). If
\(n<m\) and \(y_n=y_m\), injectivity cancels the first \(n\) iterates
and gives

\[
y_0=j^{m-n}(y_0)\in j[Y],
\]

contrary to the choice of \(y_0\). Thus \(n\mapsto y_n\) is an injection
\(\omega\hookrightarrow Y\hookrightarrow X\), contradicting Lemma 3.2. ∎

\protect\phantomsection\label{def:3-5-reversal-action}

\textbf{Definition 3.5 (Reversal action on the two types).} For each
\(i<2\), let \(\sigma(i)<2\) be the unique index satisfying

\[
L_i^*\cong L_{\sigma(i)}.
\]

The graph of \(\sigma\) is the Separation-defined subset

\[
\{(i,j)\in2\times2:L_i^*\cong L_j\};
\]

exactness makes this relation a function. Reversing twice gives

\[
\sigma^2=\operatorname{id}_{2}.
\]

Hence either \(\sigma=\operatorname{id}_{2}\) or \(\sigma=(01)\).

\protect\phantomsection\label{lem:3-6-all-selfdual-flank}

\textbf{Lemma 3.6 (Symmetric flank localization when all global types
are self-dual).} Assume every DLO on \(X\) is self-dual. Let \(P\) be a
DLO on \(X\) with an involutive reversal \(k\), and suppose

\begin{equation}\protect\phantomsection\label{eq:all-selfdual-flank}{
P=A+M+k[A],
\qquad
k[M]=M,
}\end{equation}

where \(A\) and \(k[A]\) are nonempty DLOs without endpoints. Then

\[
A\cong A^*.
\]

\emph{Proof.} On the same carrier define

\[
R=A+M+k[A]^*,
\qquad
S=A^*+M+k[A].
\]

The displayed decomposition of \(P\) makes \(M\) convex; if it has at
least two points, it is therefore internally dense. After empty blocks
are omitted, every reduced boundary is dense because the flank on the
required side has no endpoint. Hence the finite block-sum criterion
makes \(R\) and \(S\) DLOs.

The map \(k\) reverses the block order. On the reversed right flank, for
example,

\[
x<_{k[A]^*}y
\quad\Longleftrightarrow\quad
y<_{k[A]}x
\quad\Longleftrightarrow\quad
k(x)<_A k(y)
\quad\Longleftrightarrow\quad
k(y)<_{A^*}k(x).
\]

The other blocks and mixed comparisons are immediate from \(k[M]=M\) and
the fact that \(k\) reverses \(P\). Thus \(k:R\to S\) is an
anti-isomorphism, and

\[
R^*\cong S.
\]

Every DLO on \(X\) is self-dual, hence

\[
R\cong R^*\cong S.
\]

Under an isomorphism \(R\to S\), the image of the initial block \(A\) is
an initial segment of type \(A\), while the literal first block of \(S\)
is an initial segment of type \(A^*\). Corollary 2.7 gives
\(A\cong A^*\). ∎

\protect\phantomsection\label{prop:3-7-all-selfdual-dichotomy}

\textbf{Proposition 3.7 (All-self-dual rigid carrier dichotomy).} Assume
every DLO on \(X\) is self-dual and let \(L\) be a rigid DLO on \(X\).
Then

\[
X\hookrightarrow\omega
\quad\text{or}\quad
\omega\hookrightarrow X.
\]

\emph{Proof.} Choose a reversal \(h\) of \(L\). By rigidity it is unique
and involutive. \hyperref[lem:b-1-reflection-split]{Lemma B.1} gives the
convex decomposition

\[
X=H+F_h+h[H],
\]

where \(H\) and \(h[H]\) are nonempty DLOs without endpoints and \(F_h\)
is empty or a singleton. Lemma 3.6, applied with \(A=H\) and \(M=F_h\),
gives

\[
H\cong H^*.
\]

The half \(H\) is rigid because it is convex in the rigid order \(L\).
Thus the standing hypotheses of the type-count-free dyadic engine in
\hyperref[app:dyadic-engine]{Appendix B} are satisfied.
\hyperref[cor:b-8-dyadic-dichotomy]{Corollary B.8} yields

\[
X\hookrightarrow\omega
\quad\text{or}\quad
\omega\hookrightarrow X.
\]

∎

\protect\phantomsection\label{thm:3-8-residual-normal-form}

\textbf{Theorem 3.8 (Exact-two residual normal form).} If \(s(X)=2\),
then:

\begin{enumerate}
\tightlist
\item
  \(X\nhookrightarrow\omega\);
\item
  \(\omega\nhookrightarrow X\);
\item
  every DLO on \(X\) is rigid;
\item
  reversal exchanges the two types;
\item
  no DLO on \(X\) is self-dual;
\item
  after naming one representative \(L\), the spectrum is
\end{enumerate}

\begin{equation}\protect\phantomsection\label{eq:exact-two-residual}{
   \{[L],[L^*]\},
   \qquad
   L\not\cong L^*.
}\end{equation}

\emph{Proof.} The first three assertions are Lemmas 3.1--3.3. If
\(\sigma=\operatorname{id}_2\), then every DLO on \(X\) is self-dual: if
\(U\cong L_i\), then

\[
U^*\cong L_i^*\cong L_i\cong U.
\]

Choosing either rigid representative and applying Proposition 3.7 gives
\(X\hookrightarrow\omega\) or \(\omega\hookrightarrow X\), contradicting
Lemmas 3.1 and 3.2. Hence \(\sigma=(01)\). A self-dual DLO would
represent a fixed point of \(\sigma\), so none exists. The final
description follows. ∎

\protect\phantomsection\label{sec:cut-square}

\section{4. The one-point cut square}\label{the-one-point-cut-square}

Fix the representative

\[
L=(X,<)
\]

from Theorem 3.8.

\protect\phantomsection\label{def:4-1-cut-square}

\textbf{Definition 4.1 (The antipodal cut square).} For \(c\in X\), put

\[
A_c:=(-\infty,c)_L,
\qquad
B_c:=(c,\infty)_L.
\]

Define four relations on the carrier \(X\) by

\[
\mathcal C_{00}^c
=
A_c+\{c\}+B_c,
\]

\[
\mathcal C_{01}^c
=
A_c^*+\{c\}+B_c^*,
\]

\[
\mathcal C_{10}^c
=
B_c+\{c\}+A_c,
\]

and

\[
\mathcal C_{11}^c
=
B_c^*+\{c\}+A_c^*.
\]

We display the square as

\[
\begin{matrix}
\mathcal C_{00}^c & \mathcal C_{10}^c\\
\mathcal C_{01}^c & \mathcal C_{11}^c.
\end{matrix}
\]

Thus the vertical edges are \(\mathcal C_{00}^c\)--\(\mathcal C_{01}^c\)
and \(\mathcal C_{10}^c\)--\(\mathcal C_{11}^c\), while the horizontal
edges are \(\mathcal C_{00}^c\)--\(\mathcal C_{10}^c\) and
\(\mathcal C_{01}^c\)--\(\mathcal C_{11}^c\).

The order \(\mathcal C_{10}^c\) is the \textbf{cut rotation} of \(L\) at
\(c\) and is denoted

\begin{equation}\protect\phantomsection\label{eq:cut-rotation}{
\operatorname{Rot}_c(L).
}\end{equation}

\protect\phantomsection\label{lem:4-2-cut-square-actual}

\textbf{Lemma 4.2 (The cut square is actual and antipodal).} For every
\(c\in X\), all four relations in Definition 4.1 are DLOs on \(X\).
Moreover,

\begin{equation}\protect\phantomsection\label{eq:cut-square-antipodal}{
\mathcal C_{11}^c
=
(\mathcal C_{00}^c)^*,
\qquad
\mathcal C_{10}^c
=
(\mathcal C_{01}^c)^*.
}\end{equation}

\emph{Proof.} Each strict ray of a DLO is a convex DLO without
endpoints. Reversing a ray preserves this property. Inserting one
singleton between two endpointless DLO blocks satisfies the finite
block-sum criterion. Each displayed relation is defined by Separation on
\(X\times X\). Reversing the block lists gives the two identities. ∎

\protect\phantomsection\label{lem:4-3-antipodal-coloring}

\textbf{Lemma 4.3 (Antipodal two-color square).} Color the four vertices
of a \(2\times2\) square with two colors. If antipodal vertices have
opposite colors, then either both vertical edges are monochromatic or
both horizontal edges are monochromatic.

\emph{Proof.} Fix the color of the upper-left vertex. The lower-right
vertex has the other color. The upper-right vertex has one of the two
colors, and the lower-left vertex, being antipodal to it, has the other.
The two possible choices give the two conclusions. ∎

\protect\phantomsection\label{lem:4-4-selfdual-amalgamation}

\textbf{Lemma 4.4 (Two rigid self-dual DLOs and one point self-dualize
their union).} Let \(A\) and \(B\) be disjoint carriers of rigid
self-dual DLOs, and let \(c\notin A\cup B\). Then
\(A\sqcup\{c\}\sqcup B\) carries a self-dual DLO.

\emph{Proof.} Let \(r_A\) and \(r_B\) be the unique reversals. By
rigidity they are involutions. \hyperref[lem:b-1-reflection-split]{Lemma
B.1} shows that the two reflection halves of each order are nonempty
DLOs without endpoints and that each fixed-point set is empty or a
singleton. Put

\[
A^-:=\{a:a<r_A(a)\},
\qquad
A^+:=\{a:r_A(a)<a\},
\]

\[
F_A:=\{a\in A:r_A(a)=a\},
\]

and define \(B^-,B^+\) and \(F_B\) analogously. The two halves in each
pair are exchanged by the relevant reversal.

Put

\[
F:=F_A\sqcup F_B\sqcup\{c\}.
\]

Map a point of \(F_A\) to \(0\), the point \(c\) to \(1\), and a point
of \(F_B\) to \(2\). Because the three pieces are disjoint and the two
fixed-point sets are empty or singletons, this is an injection
\(F\to3\). Finite ranking \hyperref[ref:isthmus2026a]{Isthmus 2026a,
Lemma 2.8} therefore gives an enumeration

\[
F=\{z_0,\ldots,z_{t-1}\},
\qquad
1\le t\le3.
\]

Give the carrier one of the following block orders:

\[
A^-+B^-+\{z_0\}+B^++A^+
\qquad(t=1),
\]

\[
A^-+\{z_0\}+B^-+B^++\{z_1\}+A^+
\qquad(t=2),
\]

or

\[
A^-+\{z_0\}+B^-+\{z_1\}+B^++\{z_2\}+A^+
\qquad(t=3).
\]

Every displayed order is a DLO: it starts and ends with endpointless DLO
blocks, and no two singleton blocks are adjacent. Define

\begin{equation}\protect\phantomsection\label{eq:selfdual-amalgamation-map}{
h
=
\bigl(r_A\restriction(A^-\cup A^+)\bigr)
\cup
\bigl(r_B\restriction(B^-\cup B^+)\bigr)
\cup
\{(z_i,z_{t-1-i}):i<t\}.
}\end{equation}

The three domains in this union are disjoint, so \(h\) is a function.
Each restricted reversal is involutive, and \(i\mapsto t-1-i\) is an
involution on \(t\); hence

\[
h^2=\operatorname{id}_{A\sqcup\{c\}\sqcup B}.
\]

In particular, \(h\) is bijective. It reverses the entire block list and
every internal block order, so it is an order reversal of the
constructed DLO. Hence that DLO is self-dual. ∎

\protect\phantomsection\label{thm:4-5-vertical-exclusion}

\textbf{Theorem 4.5 (A vertical monochromatic cut edge is impossible).}
Under the residual normal form of Theorem 3.8, the two DLOs

\[
\mathcal C_{00}^c
\quad\text{and}\quad
\mathcal C_{01}^c
\]

are nonisomorphic for every \(c\in X\).

\emph{Proof.} Suppose

\[
f:\mathcal C_{00}^c\longrightarrow\mathcal C_{01}^c
\]

were an isomorphism.

If \(f(c)<_{\mathcal C_{01}^c}c\), then \(f(c)\in A_c\) and

\[
f[A_c]
=
(-\infty,f(c))_{\mathcal C_{01}^c}
=
\{x\in A_c:f(c)<_Lx<c\}
\subsetneq
A_c.
\]

The inclusion is proper because \(f(c)\in A_c\) but
\(f(c)\notin f[A_c]\). This is a set-theoretic statement; the target
happens to carry the reversed order on that subset. Thus
\(f\upharpoonright A_c\) is an injective nonsurjective selfmap of
\(A_c\), contradicting Lemma 3.4.

If \(c<_{\mathcal C_{01}^c}f(c)\), then \(f(c)\in B_c\) and

\[
f[B_c]
=
(f(c),\infty)_{\mathcal C_{01}^c}
=
\{x\in B_c:c<x<_L f(c)\}
\subsetneq
B_c,
\]

where the inclusion is proper because \(f(c)\in B_c\setminus f[B_c]\).
This gives the same contradiction. Therefore

\[
f(c)=c.
\]

It follows that \(f[A_c]=A_c\) and \(f[B_c]=B_c\). Since the target
order reverses the inherited order on each ray, the restrictions of
\(f\) are order reversals of the original DLOs on \(A_c\) and \(B_c\).
These ray orders are rigid by Lemma 2.8 because they are convex in the
rigid order \(L\). Lemma 4.4 therefore constructs a self-dual DLO on
\(X\), contradicting Theorem 3.8. ∎

\protect\phantomsection\label{thm:4-6-cut-rotation-type}

\textbf{Theorem 4.6 (Every cut rotation has the original type).} For
every \(c\in X\),

\begin{equation}\protect\phantomsection\label{eq:all-cut-rotations}{
\operatorname{Rot}_c(L)
\cong
L.
}\end{equation}

The isomorphism is unique.

\emph{Proof.} Color each corner \(\mathcal C_{ij}^c\) by the unique
exact-two type it represents. By Theorem 3.8, reversal exchanges the two
types. The antipodal identities in Lemma 4.2 therefore make antipodal
corners opposite in color. Lemma 4.3 gives a monochromatic direction.

The vertical direction is impossible by Theorem 4.5. Hence both
horizontal edges are monochromatic. In particular,

\[
\mathcal C_{00}^c
=
L
\cong
\mathcal C_{10}^c
=
\operatorname{Rot}_c(L).
\]

If \(f,g:L\to\operatorname{Rot}_c(L)\) are two isomorphisms, then
\(g^{-1}\circ f\) is an increasing automorphism of the rigid order
\(L\). Thus \(f=g\). ∎

\protect\phantomsection\label{sec:cut-preimage}

\section{5. The cut-preimage reversal}\label{the-cut-preimage-reversal}

Throughout this section, let \(L=(X,<)\) be a rigid DLO satisfying the
following two hypotheses:

\begin{equation}\protect\phantomsection\label{eq:hereditary-dedekind-finiteness}{
(\mathrm{HDF})
\qquad
\forall Y\subseteq X\ \forall j:Y\to Y,
\bigl(j\text{ injective}\Longrightarrow j[Y]=Y\bigr),
}\end{equation}

and

\begin{equation}\protect\phantomsection\label{eq:rotation-hypothesis}{
(\mathrm{Rot})
\qquad
\forall c\in X,
\operatorname{Rot}_c(L)\cong L.
}\end{equation}

In the exact-two application, \((\mathrm{HDF})\) is Lemma 3.4 and
\((\mathrm{Rot})\) is Theorem 4.6.

\protect\phantomsection\label{def:5-1-cut-preimage}

\textbf{Definition 5.1 (The cut-preimage map).} For a fixed \(c\in X\),
let

\[
\rho_c:
L
\longrightarrow
\operatorname{Rot}_c(L)
\]

be the unique isomorphism supplied by \((\mathrm{Rot})\) and rigidity.
Define

\begin{equation}\protect\phantomsection\label{eq:cut-preimage}{
\delta(c):=\rho_c^{-1}(c).
}\end{equation}

The graph of \(\delta\) is the Separation-defined subset of
\(X\times X\)

\begin{equation}\protect\phantomsection\label{eq:cut-preimage-graph}{
\left\{
\begin{aligned}
(c,d)\in X\times X:\;&
\exists r\in\mathcal P(X\times X)\ \text{such that}\\
&r\text{ is an order isomorphism }L\to\operatorname{Rot}_c(L),\\
&(d,c)\in r
\end{aligned}
\right\}.
}\end{equation}

Rigidity gives uniqueness of \(r\), and bijectivity of \(r\) gives
uniqueness of \(d\). No set of all maps \(\rho_c\) is formed.

\protect\phantomsection\label{lem:5-2-ray-exchange}

\textbf{Lemma 5.2 (Cut rotations exchange opposite rays).} Let
\(d=\delta(c)\). Then \(\rho_c\) restricts to order isomorphisms

\begin{equation}\protect\phantomsection\label{eq:left-to-right-ray}{
(-\infty,d)_L
\cong
(c,\infty)_L
}\end{equation}

and

\begin{equation}\protect\phantomsection\label{eq:right-to-left-ray}{
(d,\infty)_L
\cong
(-\infty,c)_L.
}\end{equation}

\emph{Proof.} In the rotated order

\[
\operatorname{Rot}_c(L)
=
(c,\infty)_L+\{c\}+(-\infty,c)_L,
\]

the strict initial segment below \(c\) is the original right ray and the
strict final segment above \(c\) is the original left ray. Since
\(\rho_c(d)=c\), an order isomorphism sends the two rays at \(d\) to
these two segments. ∎

\protect\phantomsection\label{lem:5-3-delta-injective}

\textbf{Lemma 5.3 (The cut-preimage map is injective).} The function
\(\delta:X\to X\) is injective.

\emph{Proof.} Suppose

\[
\delta(c)=\delta(e)=d
\]

with \(c<e\). Put

\[
A:=(-\infty,d)_L,
\qquad
T_c:=(c,\infty)_L,
\qquad
T_e:=(e,\infty)_L.
\]

Lemma 5.2 gives isomorphisms

\[
\rho_c\upharpoonright A:A\to T_c,
\qquad
\rho_e\upharpoonright A:A\to T_e.
\]

Therefore

\[
T_c
\xrightarrow{\ (\rho_c\upharpoonright A)^{-1}\ }
A
\xrightarrow{\ \rho_e\upharpoonright A\ }
T_e
\hookrightarrow
T_c
\]

is an injective selfmap of \(T_c\) with range exactly \(T_e\). The
inclusion is proper because

\[
e\in T_c\setminus T_e.
\]

This contradicts \((\mathrm{HDF})\). The case \(e<c\) is symmetric. ∎

\protect\phantomsection\label{lem:5-4-delta-decreasing}

\textbf{Lemma 5.4 (The cut-preimage map is strictly decreasing).} If
\(c<e\), then

\begin{equation}\protect\phantomsection\label{eq:delta-decreasing}{
\delta(e)<\delta(c).
}\end{equation}

\emph{Proof.} Suppose instead that

\[
\delta(c)<\delta(e).
\]

Put

\[
I_c:=(-\infty,\delta(c))_L,
\qquad
I_e:=(-\infty,\delta(e))_L.
\]

Then

\[
I_c\subsetneq I_e,
\]

with

\[
\delta(c)\in I_e\setminus I_c.
\]

Meanwhile,

\[
(e,\infty)_L
\subsetneq
(c,\infty)_L.
\]

Using Lemma 5.2, form the composite

\begin{equation}\protect\phantomsection\label{eq:proper-ray-selfmap}{
I_e
\xrightarrow{\ \rho_e\ }
(e,\infty)_L
\hookrightarrow
(c,\infty)_L
\xrightarrow{\ \rho_c^{-1}\ }
I_c
\hookrightarrow
I_e.
}\end{equation}

This is an injective selfmap of \(I_e\), and

\begin{equation}\protect\phantomsection\label{eq:proper-ray-selfmap-range}{
\operatorname{ran}(\Psi_{c,e})
=
\rho_c^{-1}[(e,\infty)_L]
\subseteq
I_c
\subsetneq
I_e,
}\end{equation}

where \(\Psi_{c,e}\) denotes the displayed composite. Thus it is not
surjective, contradicting \((\mathrm{HDF})\). Hence
\(\delta(c)\not<\delta(e)\). Lemma 5.3 excludes equality, so
\(\delta(e)<\delta(c)\). ∎

\protect\phantomsection\label{cor:5-5-delta-bijective}

\textbf{Corollary 5.5 (The cut-preimage map is a decreasing bijection).}
The function \(\delta:X\to X\) is bijective and strictly decreasing.

\emph{Proof.} It is injective by Lemma 5.3. Hypothesis
\((\mathrm{HDF})\), applied to \(Y=X\), makes every injective selfmap of
\(X\) surjective. Strict decrease is Lemma 5.4. ∎

\protect\phantomsection\label{thm:5-6-cut-rotation-selfduality}

\textbf{Theorem 5.6 (Cut-rotation self-duality theorem).} Let
\(L=(X,<)\) be a rigid DLO satisfying \((\mathrm{HDF})\) and
\((\mathrm{Rot})\). Then

\[
L\cong L^*.
\]

\emph{Proof.} Rigidity makes every cut-rotation isomorphism unique, so
Definition 5.1 gives the cut-preimage map. Lemmas 5.2--5.4 and Corollary
5.5 show that \(\delta\) is a decreasing bijection of \(X\). Hence, for
all \(x,y\in X\),

\begin{equation}\protect\phantomsection\label{eq:delta-anti-isomorphism}{
x<_L y
\quad\Longleftrightarrow\quad
\delta(y)<_L\delta(x)
\quad\Longleftrightarrow\quad
\delta(x)<_{L^*}\delta(y).
}\end{equation}

For the reverse implication in the first equivalence, suppose
\(\delta(y)<_L\delta(x)\). If \(y<_Lx\), strict decrease gives
\(\delta(x)<_L\delta(y)\), while \(x=y\) contradicts the assumed strict
inequality. Trichotomy therefore gives \(x<_Ly\). Hence \(\delta\) is an
order isomorphism from \(L\) to \(L^*\). ∎

\protect\phantomsection\label{sec:main-theorem}

\section{6. Exact-two exclusion}\label{exact-two-exclusion}

\protect\phantomsection\label{thm:6-1-main}

\textbf{Theorem 6.1 (No exact-two DLO spectrum).} \(Z_{\mathrm{sep}}\)
proves

\begin{equation}\protect\phantomsection\label{eq:main-theorem}{
\neg\exists X\,(s(X)=2).
}\end{equation}

\emph{Proof.} Assume \(s(X)=2\). Theorem 3.8 gives a rigid
representative \(L\) satisfying

\[
L\not\cong L^*,
\]

Lemma 3.4 gives \((\mathrm{HDF})\), and Theorem 4.6 gives
\((\mathrm{Rot})\). Theorem 5.6 therefore gives \(L\cong L^*\), a
contradiction. ∎

\protect\phantomsection\label{cor:6-2-two-three-gap}

\textbf{Corollary 6.2 (No exact two or exact three).} One has

\begin{equation}\protect\phantomsection\label{eq:two-three-gap}{
Z_{\mathrm{sep}}
\vdash
\neg\exists X\,
\bigl(
\operatorname{Spec}_{=2}(X)
\vee
\operatorname{Spec}_{=3}(X)
\bigr).
}\end{equation}

\emph{Proof.} The exact-two case is Theorem 6.1. The exact-three case is
\hyperref[ref:isthmus2026a]{Isthmus 2026a, Theorem 6.3}. ∎

\protect\phantomsection\label{cor:6-3-zf}

\textbf{Corollary 6.3 (ZF consequence).} ZF proves that no set carries
exactly two DLO isomorphism types.

\emph{Proof.} ZF contains every axiom of \(Z_{\mathrm{sep}}\). ∎

\protect\phantomsection\label{sec:scope}

\section{7. Scope and further
questions}\label{scope-and-further-questions}

For each standard finite \(n\ge2\), let \((\mathrm{FS})_n\) denote the
metatheoretic assertion

\[
Z_{\mathrm{sep}}
\vdash
\neg\exists X\,
\operatorname{Spec}_{=n}(X).
\]

The present paper and its exact-three companion prove
\((\mathrm{FS})_2\) and \((\mathrm{FS})_3\). They do not prove
\((\mathrm{FS})_n\) for \(n\ge4\).

The exact-two mechanism is specific to a two-color antipodal square.
Once reversal is known to exchange the two types, opposite cut-square
corners have opposite colors, and one of the two edge directions is
forced to be monochromatic. Vertical monochromaticity is incompatible
with the absence of self-dual types, so every cut rotation has the
original type. For larger spectra, a four-corner cut square need not
force any same-type edge, and the argument stops before the cut-preimage
map can be defined globally.

Theorem 5.6 is independent of finite spectrum language. It says that a
rigid DLO on a hereditarily Dedekind-finite carrier cannot be isomorphic
to all of its one-point cut rotations unless it is self-dual. This
cut-rotation rigidity principle may be useful in other choiceless
classification problems.

The present methods do not decide whether any standard finite spectrum
of size at least four can occur, or which such spectra are consistent.
The countable-benchmark theorem gives only the lower bound \(s(X)\ge4\)
on a non-subcountable Dedekind-infinite carrier; that bound does not
exclude exact four or any larger finite value.

\protect\phantomsection\label{app:weak-audit}

\section{Appendix A. Weak-theory audit of the new
construction}\label{appendix-a.-weak-theory-audit-of-the-new-construction}

\protect\phantomsection\label{prop:a-1-cut-relations}

\textbf{Proposition A.1 (Uniform formation of the cut rotations).} For
each \(c\in X\), the graph of \(\operatorname{Rot}_c(L)\) is a set
obtained by Separation on \(X\times X\).

\emph{Proof.} For \(x,y\in X\), the relation
\(x<_{\operatorname{Rot}_c(L)}y\) is the finite disjunction saying that
either:

\begin{enumerate}
\tightlist
\item
  \(c<x<y\) in \(L\);
\item
  \(c<x\) and \(y=c\);
\item
  \(c<x\) and \(y<c\);
\item
  \(x=c\) and \(y<c\);
\item
  \(x<y<c\) in \(L\).
\end{enumerate}

This formula defines the block order

\[
(c,\infty)_L+\{c\}+(-\infty,c)_L
\]

inside the fixed set \(X\times X\). ∎

\protect\phantomsection\label{prop:a-2-cut-table}

\textbf{Proposition A.2 (The finite type table requires no Choice).} For
a fixed \(c\), assigning to each of the four cut-square relations its
unique exact-two representative is a finite set construction in
\(Z_{\mathrm{sep}}\).

\emph{Proof.} The four relations form a set by Pairing and Union. For
each relation, exactness gives a unique index in \(2\). The graph of the
assignment is the Separation subset of the finite product \(4\times2\)
satisfying the internal isomorphism predicate. No choice from an
infinite family is involved. ∎

\protect\phantomsection\label{prop:a-3-delta-graph}

\textbf{Proposition A.3 (The cut-preimage graph uses no Replacement).}
The graph in Equation
\hyperref[eq:cut-preimage-graph]{\eqref{eq:cut-preimage-graph}} is a set
in \(Z_{\mathrm{sep}}\).

\emph{Proof.} Every candidate is a pair in the fixed set \(X\times X\),
and every candidate isomorphism is an element of the fixed power set
\(\mathcal P(X\times X)\). The formula in Equation
\hyperref[eq:cut-preimage-graph]{\eqref{eq:cut-preimage-graph}}
therefore defines the graph by Separation. Uniqueness turns it into a
function. At no point is the range

\[
\{\rho_c:c\in X\}
\]

formed. ∎

\protect\phantomsection\label{prop:a-4-composites}

\textbf{Proposition A.4 (The ray composites are internal maps).} Every
restriction, inverse, inclusion, and finite composite used in Lemmas 5.3
and 5.4 is an internal set.

\emph{Proof.} Restrictions and inclusions are Separation-defined subsets
of \(X\times X\). The inverse of a bijection is obtained by reversing
ordered pairs. The graph of a composite of two maps is defined by an
existential formula on the appropriate fixed Cartesian product. Only
finitely many maps are composed in each argument, so Pairing, Union,
products, and Separation suffice. ∎

\protect\phantomsection\label{app:dyadic-engine}

\section{Appendix B. A type-count-free dyadic reflection
engine}\label{appendix-b.-a-type-count-free-dyadic-reflection-engine}

This appendix isolates the reflection-tree mechanism needed in
Proposition 3.7. It is parameterized by self-duality inputs rather than
by an exact finite spectrum. In addition to the standing infrastructure
recorded in \hyperref[sec:formal-setting]{§2}, it uses Lemma 2.8 and
Lemma 3.6. Its remaining imported inputs are elementary foundational
lemmas from the exact-three companion: finite ranking, finite unions of
subcountable sets, explicit binary-word coding, and internally finite
block sums \hyperref[ref:isthmus2026a]{Isthmus 2026a, Lemmas 2.8,
3.4(1), B.3, and Proposition B.5}. No result whose statement assumes
\(s(X)=3\) is used below.

To keep word sets distinct from finite ordinals, write

\[
\operatorname{Bin}_{<\omega}
:=
2^{<\omega},
\qquad
\operatorname{Bin}_n
:=
\{s\in\operatorname{Bin}_{<\omega}:|s|=n\},
\]

\[
\operatorname{Bin}_{<n}
:=
\{s\in\operatorname{Bin}_{<\omega}:|s|<n\},
\qquad
\operatorname{Bin}_{\le n}
:=
\{s\in\operatorname{Bin}_{<\omega}:|s|\le n\},
\]

and let

\[
N_n:=2^n.
\]

Here the right-hand side is the finite ordinal obtained by
exponentiation, so \(N_n\in\omega\). We use the explicit bijections

\[
v_n:\operatorname{Bin}_n\longrightarrow N_n
\]

and

\[
c_n:\operatorname{Bin}_{<n}\longrightarrow N_n-1
\]

from \hyperref[ref:isthmus2026a]{Isthmus 2026a, Lemma B.3}.

\protect\phantomsection\label{lem:b-1-reflection-split}

\textbf{Lemma B.1 (Geometry of one reflection split).} Let \(I\) be a
nonempty rigid self-dual DLO and let \(r\) be its unique reversal. Put

\[
I_0:=\{x\in I:x<r(x)\},
\]

\[
I_1:=\{x\in I:r(x)<x\},
\]

and

\[
F_r:=\{x\in I:r(x)=x\}.
\]

Then:

\begin{enumerate}
\tightlist
\item
  \(r\) is involutive;
\item
  \(F_r\) is empty or a singleton;
\item
  \(I_0\) and \(I_1\) are nonempty proper convex DLOs without endpoints;
\item
  \(I=I_0\sqcup F_r\sqcup I_1\);
\item
  \(r\) anti-isomorphically exchanges \(I_0\) and \(I_1\).
\end{enumerate}

\emph{Proof.} Lemma 2.8 gives involutivity. An order reversal has at
most one fixed point: if \(x<y\) were both fixed, reversal would give
\(y=r(y)<r(x)=x\). The two strict-inequality sets are respectively
initial and final, are convex, and are exchanged by \(r\).

They are nonempty because \(r\) cannot fix every point of a nontrivial
strict order, by the same two-point argument. To show that \(I_0\) has
no greatest point, take \(x\in I_0\) and choose

\[
x<z_1<z_2<r(x).
\]

At most one of \(z_1,z_2\) is fixed by \(r\), so choose a nonfixed
\(z\).

\begin{itemize}
\tightlist
\item
  If \(z<r(z)\), then \(z\in I_0\) and \(x<z\).
\item
  If \(r(z)<z\), reversing \(z<r(x)\) gives \(x<r(z)\), and
  \(r(z)\in I_0\).
\end{itemize}

The remaining endpoint statements are symmetric or follow from
initiality and finality. Density is inherited from convexity, and
trichotomy gives the displayed decomposition. ∎

\protect\phantomsection\label{lem:b-2-singleton-placement}

\textbf{Lemma B.2 (Sharp singleton placement).} Let \(q\ge1\). Let
\(C_0,\ldots,C_{q-1}\) be nonempty DLO blocks without endpoints and let
\(z_0,\ldots,z_{t-1}\) be singleton blocks. If \(t\le q\), the order

\begin{equation}\protect\phantomsection\label{eq:typefree-singleton-placement}{
z_0+C_0+z_1+C_1+\cdots+z_{t-1}+C_{t-1}
+C_t+\cdots+C_{q-1}
}\end{equation}

with the initial alternating part omitted when \(t=0\) and the final
tail omitted when \(t=q\), has no greatest point and contains no
adjacent singleton blocks.

If \(t=q+1\), no ordering of the \(q\) DLO blocks and the \(t\)
singleton blocks can simultaneously avoid adjacent singletons and avoid
a singleton final block.

\emph{Proof.} The displayed arrangement proves the first assertion. For
the second, \(q\) nonsingleton blocks provide at most \(q\) separators
between singleton blocks. Accommodating \(q+1\) singleton blocks without
adjacency forces the alternating pattern to begin and end with a
singleton. ∎

For the rest of the appendix, fix a rigid DLO \(L=(X,<)\) with an
involutive reversal \(h\), and write

\begin{equation}\protect\phantomsection\label{eq:typefree-root-decomposition}{
X=H+F_h+h[H]
}\end{equation}

as in Lemma B.1. Assume that \(H\) is self-dual and that every DLO on
\(X\) is self-dual.

\protect\phantomsection\label{lem:b-3-child-localization}

\textbf{Lemma B.3 (All-self-dual finite localization).} Let
\(n<\omega\). Suppose that:

\begin{enumerate}
\tightlist
\item
  \(H\) is partitioned into cells \((I_s)_{s\in\operatorname{Bin}_n}\)
  and a center set \(P_{<n}\);
\item
  every \(I_s\) is a nonempty convex rigid self-dual DLO;
\item
  there is an injective tag map
\end{enumerate}

\[
   \tau_n:P_{<n}\longrightarrow\operatorname{Bin}_{<n}.
\]

For each \(s\in\operatorname{Bin}_n\), let \(r_s\) be the unique
reversal of \(I_s\) and put

\[
F_s:=\{x\in I_s:r_s(x)=x\}.
\]

Split \(I_s\) by \(r_s\), and, when \(F_s\) is nonempty, tag its unique
point by \(s\in\operatorname{Bin}_n\). Then every resulting
level-\((n+1)\) child cell \(D\) satisfies

\begin{equation}\protect\phantomsection\label{eq:typefree-child-selfdual}{
D\cong D^*.
}\end{equation}

\emph{Proof.} The child cells are indexed by
\(\operatorname{Bin}_{n+1}\). Old center tags lie in
\(\operatorname{Bin}_{<n}\) and new center tags lie in
\(\operatorname{Bin}_n\). Old centers lie outside every level-\(n\)
cell, while a new center lies inside its parent cell; new centers from
distinct parents lie in disjoint cells. Hence the combined tag
assignment is injective into \(\operatorname{Bin}_{<n+1}\).

Let \(d\in\operatorname{Bin}_{n+1}\) index the target cell \(D\), put
\(a=v_{n+1}(d)\), and define

\[
\rho(s)
=
\begin{cases}
 v_{n+1}(s),&v_{n+1}(s)<a,\\
 v_{n+1}(s)-1,&a<v_{n+1}(s).
\end{cases}
\]

Then

\begin{equation}\protect\phantomsection\label{eq:typefree-remaining-cell-ranking}{
\rho:
\{s\in\operatorname{Bin}_{n+1}:s\ne d\}
\longrightarrow
q,
\qquad
q:=N_{n+1}-1,
}\end{equation}

is a bijection. Thus the other child cells have a canonical enumeration

\[
C_0,\ldots,C_{q-1}.
\]

Composing the center tags with
\(c_{n+1}:\operatorname{Bin}_{<n+1}\to q\) injects the total center set
into \(q\). Finite ranking \hyperref[ref:isthmus2026a]{Isthmus 2026a,
Lemma 2.8} therefore gives some \(t\le q\) and an enumeration

\[
z_0,\ldots,z_{t-1}
\]

of all accumulated center points. Lemma B.2 gives a block order \(W_D\)
on \(H\setminus D\) in which every center is followed by a DLO cell and
the final block is a DLO:

\begin{equation}\protect\phantomsection\label{eq:typefree-localized-complement}{
W_D
=
z_0+C_0+\cdots+z_{t-1}+C_{t-1}
+C_t+\cdots+C_{q-1}.
}\end{equation}

Here the initial alternating part is omitted when \(t=0\), and the final
tail is omitted when \(t=q\). As sets,

\begin{equation}\protect\phantomsection\label{eq:typefree-localized-carrier-partition}{
H=D\sqcup W_D,
\qquad
X=D\sqcup W_D\sqcup F_h\sqcup h[W_D]\sqcup h[D].
}\end{equation}

The combined tag graph is the union of the old tag map with the
Separation-defined graph

\[
\{(x,s)\in H\times\operatorname{Bin}_n:x\in F_s\},
\]

and the enumeration \(i\mapsto C_i\) is the Separation-defined subset of
\(q\times\mathcal P(H)\) determined by \(\rho\). Here and below, \(W_D\)
denotes both the displayed ordered block sum and its carrier
\(H\setminus D\).

Let

\[
\mathfrak B_D\subseteq(q+t)\times\mathcal P(H)
\]

be the graph of the displayed block family, and let

\[
\mathfrak R_D\subseteq(q+t)\times\mathcal P(H\times H)
\]

be the graph assigning to each singleton block the empty relation and to
each cell block \(C_i\) its order inherited from \(L\). Both graphs are
formed by Separation from the two finite enumerations. Applying
\hyperref[ref:isthmus2026a]{Isthmus 2026a, Proposition B.5} to
\((\mathfrak B_D,\mathfrak R_D)\) yields the graph of the internally
finite block-sum order \(W_D\).

Give \(D\) and \(h[D]\) their orders inherited from \(L\). Give
\(h[W_D]\) the mirror order

\begin{equation}\protect\phantomsection\label{eq:typefree-mirror-order}{
u<_{h[W_D]}v
\quad\Longleftrightarrow\quad
h(v)<_{W_D}h(u).
}\end{equation}

Put

\[
M_D:=W_D+F_h+h[W_D]
\]

and

\begin{equation}\protect\phantomsection\label{eq:typefree-localized-order}{
L_D:=D+M_D+h[D].
}\end{equation}

After empty blocks are omitted, the following list covers every
boundary.

{\def\LTcaptype{none} % do not increment counter
\begin{longtable}[]{@{}
  >{\raggedright\arraybackslash}p{(\linewidth - 4\tabcolsep) * \real{0.3333}}
  >{\raggedright\arraybackslash}p{(\linewidth - 4\tabcolsep) * \real{0.3333}}
  >{\raggedright\arraybackslash}p{(\linewidth - 4\tabcolsep) * \real{0.3333}}@{}}
\toprule\noalign{}
\begin{minipage}[b]{\linewidth}\raggedright
Boundary
\end{minipage} & \begin{minipage}[b]{\linewidth}\raggedright
Range
\end{minipage} & \begin{minipage}[b]{\linewidth}\raggedright
Reason for density
\end{minipage} \\
\midrule\noalign{}
\endhead
\bottomrule\noalign{}
\endlastfoot
\(D\mid z_0\) & \(t>0\) & \(D\) has no greatest point \\
\(D\mid C_0\) & \(t=0\) & \(D\) has no greatest point \\
\(z_i\mid C_i\) & \(0\le i<t\) & \(C_i\) has no least point \\
\(C_i\mid z_{i+1}\) & \(0\le i<t-1\) & \(C_i\) has no greatest point \\
\(C_j\mid C_{j+1}\) & \(0\le j<q-1\) if \(t=0\); \(t-1\le j<q-1\) if
\(0<t<q\); none if \(t=q\) & \(C_j\) has no greatest point \\
\(W_D\mid F_h\) & \(F_h\ne\varnothing\) & \(W_D\) has no greatest
point \\
\(F_h\mid h[W_D]\) & \(F_h\ne\varnothing\) & \(h[W_D]\) has no least
point \\
\(W_D\mid h[W_D]\) & \(F_h=\varnothing\) & \(W_D\) has no greatest point
and \(h[W_D]\) has no least point \\
internal mirror boundaries & all corresponding indices & duals of the
preceding internal cases \\
\(h[W_D]\mid h[D]\) & always & \(h[D]\) has no least point \\
\end{longtable}
}

The first block \(D\) has no least point and the last block \(h[D]\) has
no greatest point. Hence the finite block-sum criterion makes \(L_D\) a
DLO on the same carrier \(X\).

The map \(h\) reverses \(L_D\), exchanges \(D\) with \(h[D]\), and
satisfies

\[
h[M_D]=M_D.
\]

Every DLO on \(X\) is self-dual, so Lemma 3.6 applies to

\[
L_D=D+M_D+h[D]
\]

and gives \(D\cong D^*\). ∎

\protect\phantomsection\label{def:b-4-level-state}

\textbf{Definition B.4 (Complete dyadic level states).} Put

\begin{equation}\protect\phantomsection\label{eq:typefree-tree-record-space}{
U_H
:=
\operatorname{Bin}_{<\omega}
\times
\mathcal P(H)
\times
\mathcal P(H\times H)
\times
\mathcal P(H).
}\end{equation}

A record \((s,I,r,F)\in U_H\) consists of a node, a cell, a reversal
graph, and its fixed-point set.

For \(n<\omega\), a subset \(\mathcal S\subseteq U_H\) is a
\textbf{valid level-\(n\) state} if:

\begin{enumerate}
\tightlist
\item
  for every \(s\in\operatorname{Bin}_{\le n}\), exactly one record
  \((s,I_s,r_s,F_s)\) belongs to \(\mathcal S\), and no record with
  longer node belongs to \(\mathcal S\);
\item
  \(I_\varnothing=H\);
\item
  \(r_s\) is the unique reversal of \(I_s\), and
\end{enumerate}

\[
   F_s=\{x\in I_s:r_s(x)=x\};
\]

\begin{enumerate}
\setcounter{enumi}{3}
\tightlist
\item
  for \(|s|<n\),
\end{enumerate}

\[
   I_{s0}=\{x\in I_s:x<r_s(x)\},
\]

\[
   I_{s1}=\{x\in I_s:r_s(x)<x\};
\]

\begin{enumerate}
\setcounter{enumi}{4}
\tightlist
\item
  every \(I_s\) is nonempty, convex in \(L\), rigid, self-dual, and a
  DLO;
\item
  for every \(k\le n\),
\end{enumerate}

\begin{equation}\protect\phantomsection\label{eq:typefree-level-partition}{
   H
   =
   \left(\bigsqcup_{|s|=k}I_s\right)
   \sqcup
   P_{<k},
}\end{equation}

where

\[
   P_{<k}
   :=
   \{x\in H:\exists s\ (|s|<k\text{ and }x\in F_s)\},
\]

and there is an injection

\[
   P_{<k}\longrightarrow N_k-1.
\]

Write \(\operatorname{Valid}_H(n,\mathcal S)\) for the conjunction of
these clauses.

\protect\phantomsection\label{thm:b-5-tree-actualization}

\textbf{Theorem B.5 (Type-count-free dyadic actualization).} For every
\(n<\omega\), there is a unique valid level-\(n\) state
\(\mathcal S_n\). The states are coherent under restriction, and the
complete tree

\begin{equation}\protect\phantomsection\label{eq:typefree-complete-tree}{
\mathcal T_H
:=
\left\{
 u\in U_H:
 \exists n\in\omega\,
 \exists\mathcal S\in\mathcal P(U_H)\,
 \bigl(
  \operatorname{Valid}_H(n,\mathcal S)
  \wedge
  u\in\mathcal S
 \bigr)
\right\}
}\end{equation}

is a set in \(Z_{\mathrm{sep}}\).

\emph{Proof.} For \(n=0\), the standing hypotheses make \(H\) rigid and
self-dual. Its reversal and fixed-point set are unique, so
\(\mathcal S_0\) exists uniquely.

Assume \(\mathcal S_n\) exists uniquely. Every old center has a unique
node tag: a point in \(F_s\) belongs to neither child of \(s\) and hence
to no descendant cell, while cells at incomparable nodes lie in disjoint
branches. Since each \(F_s\) is empty or a singleton, the map

\begin{equation}\protect\phantomsection\label{eq:typefree-old-center-tags}{
P_{<n}\longrightarrow\operatorname{Bin}_{<n},
\qquad
x\longmapsto\text{the unique }s\text{ with }x\in F_s,
}\end{equation}

is injective.

Lemma B.1 determines the two geometric children of every level-\(n\)
cell. Lemma B.3, applied with the preceding tag map, makes every child
self-dual. Each child is convex in \(L\) and therefore rigid by Lemma
2.8. Its reversal and fixed-point set are consequently unique.

For \(u\in U_H\), let \(\operatorname{Old}_{\mathcal S_n}(u)\) mean
\(u\in\mathcal S_n\). Let \(\operatorname{New}_{\mathcal S_n}(u)\) mean
that

\[
u=(s,I,r,F),
\qquad
s=t^\frown\varepsilon,
\qquad
|t|=n,
\qquad
\varepsilon\in2,
\]

there is a parent record \((t,I_t,r_t,F_t)\in\mathcal S_n\), and:

\begin{itemize}
\tightlist
\item
  if \(\varepsilon=0\), then \(I=\{x\in I_t:x<r_t(x)\}\);
\item
  if \(\varepsilon=1\), then \(I=\{x\in I_t:r_t(x)<x\}\);
\item
  \(r\) is the unique reversal of \(I\);
\item
  \(F=\{x\in I:r(x)=x\}\).
\end{itemize}

Put

\begin{equation}\protect\phantomsection\label{eq:typefree-successor-state}{
\mathcal S_{n+1}
=
\{u\in U_H:
\operatorname{Old}_{\mathcal S_n}(u)
\vee
\operatorname{New}_{\mathcal S_n}(u)\}.
}\end{equation}

This is a Separation-defined subset of the fixed set \(U_H\).

Tag old centers by their old nodes and a possible new center of \(I_s\)
by \(s\in\operatorname{Bin}_n\). This gives an injection

\[
P_{<n+1}\longrightarrow\operatorname{Bin}_{<n+1}\xrightarrow{\ c_{n+1}\ }N_{n+1}-1.
\]

Clauses 1--5 of Definition B.4 hold by construction. For \(k\le n\), the
partition and center bound are inherited from \(\mathcal S_n\). For
\(k=n+1\), Lemma B.1 partitions every parent into its two children and
possible fixed point, while the displayed tag injection gives the
required bound. Hence clause 6 holds as well. The geometric children are
uniquely determined, and rigidity makes their reversals and fixed-point
sets unique, so \(\mathcal S_{n+1}\) is the unique valid level-\((n+1)\)
state.

If \(m>n\), the restriction of any valid level-\(m\) state to nodes of
length at most \(n\) is a valid level-\(n\) state and hence equals
\(\mathcal S_n\). Thus the states are coherent. Finally,
\(\operatorname{Valid}_H(n,\mathcal S)\) is a formula on the fixed set
\(\omega\times\mathcal P(U_H)\), and the displayed definition of
\(\mathcal T_H\) is Separation on \(U_H\); no Replacement-generated
range is formed. ∎

\protect\phantomsection\label{lem:b-6-center-countability}

\textbf{Lemma B.6 (The center set is at most countable).} Define

\begin{equation}\protect\phantomsection\label{eq:typefree-center-set}{
Z_H
:=
\{x\in H:\exists(s,I,r,F)\in\mathcal T_H\ (x\in F)\}.
}\end{equation}

Then

\begin{equation}\protect\phantomsection\label{eq:typefree-center-countable}{
Z_H\hookrightarrow\operatorname{Bin}_{<\omega}\hookrightarrow\omega.
}\end{equation}

\emph{Proof.} Every center belongs to \(F_s\) for a unique node \(s\):
centers are excluded from all descendant cells, and incomparable nodes
lie in disjoint cells. Two centers with the same tag lie in the
empty-or-singleton set \(F_s\) and are equal. Thus \(x\mapsto s\) is an
injection obtained by Separation from
\(H\times\operatorname{Bin}_{<\omega}\). The explicit finite-word code
of \hyperref[ref:isthmus2026a]{Isthmus 2026a, Lemma B.3} gives
\(\operatorname{Bin}_{<\omega}\hookrightarrow\omega\). ∎

\protect\phantomsection\label{thm:b-7-off-center}

\textbf{Theorem B.7 (An off-center point generates an injective
sequence).} Let \(x\in H\setminus Z_H\). For every \(n\), let \(s_n\) be
the unique word of length \(n\) whose cell contains \(x\), and put

\begin{equation}\protect\phantomsection\label{eq:typefree-off-center-sequence}{
y_n:=r_{s_n}(x).
}\end{equation}

Then \(n\mapsto y_n\) is an injection \(\omega\hookrightarrow X\).

\emph{Proof.} The level partition and \(x\notin Z_H\) give the unique
cell \(I_{s_n}\). Since \(x\) is not fixed by \(r_{s_n}\), the points
\(x\) and \(y_n\) lie in opposite children. The next chosen cell
\(I_{s_{n+1}}\) is the child containing \(x\), so

\[
y_n\notin I_{s_{n+1}}.
\]

If \(m>n\), then

\[
y_m\in I_{s_m}\subseteq I_{s_{n+1}},
\]

and hence \(y_m\ne y_n\).

The graph is the Separation-defined subset of \(\omega\times X\)

\begin{equation}\protect\phantomsection\label{eq:typefree-off-center-graph}{
G_x
:=
\left\{
(n,y)\in\omega\times X:
\exists s,I,r,F\,
\left[
\begin{aligned}
 &(s,I,r,F)\in\mathcal T_H,\\
 &|s|=n,\ x\in I,\ (x,y)\in r
\end{aligned}
\right]
\right\}.
}\end{equation}

Coherence and record uniqueness make \(G_x\) a function, and the
preceding argument makes it injective. ∎

\protect\phantomsection\label{cor:b-8-dyadic-dichotomy}

\textbf{Corollary B.8 (Type-count-free rigid carrier dichotomy).} Under
the standing hypotheses of this appendix,

\begin{equation}\protect\phantomsection\label{eq:typefree-rigid-dichotomy}{
X\hookrightarrow\omega
\quad\text{or}\quad
\omega\hookrightarrow X.
}\end{equation}

\emph{Proof.} If \(H\hookrightarrow\omega\), then
\(h[H]\hookrightarrow\omega\) and \(F_h\) is empty or a singleton. The
finite-union coding of \hyperref[ref:isthmus2026a]{Isthmus 2026a, Lemma
3.4(1)} gives

\[
X=H+F_h+h[H]\hookrightarrow\omega.
\]

If \(H\nhookrightarrow\omega\), Lemma B.6 implies \(H\ne Z_H\). Choose
one point in \(H\setminus Z_H\) and apply Theorem B.7. ∎

\protect\phantomsection\label{sec:ai-disclosure}

\section*{AI-use disclosure}\label{ai-use-disclosure}

AI systems were used extensively in an iterative, human-directed
workflow. OpenAI ChatGPT was used for exploratory proof development,
counterexample search, drafting, and virtual review; Anthropic Claude
Code was used for Lean formalization, proof repair, and validation; and
OpenAI Codex was used for mathematical and formalization audits, release
review, and editorial checks. AI assistance was also used to translate
and polish English prose. The author set the research objectives,
decided which candidate arguments to retain or reject, reviewed and
approved the final manuscript and formalization, and takes
responsibility for all mathematical claims and any remaining errors. AI
systems are tools, not authors.

\protect\phantomsection\label{sec:references}

\section{References}\label{references}

\protect\phantomsection\label{ref:ervin2017}

\begin{itemize}
\tightlist
\item
  \textbf{Ervin 2017.} Garrett Ervin. ``Every Linear Order Isomorphic to
  Its Cube Is Isomorphic to Its Square.'' \emph{Advances in Mathematics}
  313 (2017): 237--281. DOI: 10.1016/j.aim.2017.04.010.
\end{itemize}

\protect\phantomsection\label{ref:howardrubin1998}

\begin{itemize}
\tightlist
\item
  \textbf{Howard--Rubin 1998.} Paul Howard and Jean E. Rubin.
  \emph{Consequences of the Axiom of Choice}. Mathematical Surveys and
  Monographs 59. American Mathematical Society, 1998.
\end{itemize}

\protect\phantomsection\label{ref:isthmus2026a}

\begin{itemize}
\tightlist
\item
  \textbf{Isthmus 2026a.} Lior Isthmus. ``No Set Carries Exactly Three
  Dense Linear Orders without Endpoints: A Proof in a Weak Zermelo
  Theory without Choice or Replacement.'' Preprint, version
  \texttt{paper-rc6}, 2026. DOI:
  \href{https://doi.org/10.5281/zenodo.21735262}{10.5281/zenodo.21735262}.
  \href{https://github.com/lioristhmus/no-exactly-three-dlo}{GitHub}.
\end{itemize}

\protect\phantomsection\label{ref:isthmus2026b}

\begin{itemize}
\tightlist
\item
  \textbf{Isthmus 2026b.} Lior Isthmus. \emph{No Set Carries Exactly
  Three Dense Linear Orders without Endpoints: Lean 4 Formalization}.
  Software release.
  \href{https://github.com/lioristhmus/no-exactly-three-dlo-lean}{GitHub}.
  DOI:
  \href{https://doi.org/10.5281/zenodo.21729423}{10.5281/zenodo.21729423}.
\end{itemize}

\protect\phantomsection\label{ref:isthmus2026c}

\begin{itemize}
\tightlist
\item
  \textbf{Isthmus 2026c.} Lior Isthmus. \emph{No Set Carries Exactly Two
  Dense Linear Orders without Endpoints: Lean 4 Formalization}. Software
  release candidate.
  \href{https://github.com/lioristhmus/no-exactly-two-dlo-lean}{GitHub}.
\end{itemize}

\protect\phantomsection\label{ref:jech1973}

\begin{itemize}
\tightlist
\item
  \textbf{Jech 1973.} Thomas J. Jech. \emph{The Axiom of Choice}.
  Studies in Logic and the Foundations of Mathematics 75. North-Holland,
  1973.
\end{itemize}

\protect\phantomsection\label{ref:marker2002}

\begin{itemize}
\tightlist
\item
  \textbf{Marker 2002.} David Marker. \emph{Model Theory: An
  Introduction}. Graduate Texts in Mathematics 217. Springer, 2002.
\end{itemize}

\protect\phantomsection\label{ref:shelah2009}

\begin{itemize}
\tightlist
\item
  \textbf{Shelah 2009.} Saharon Shelah. ``Model Theory without Choice?
  Categoricity.'' \emph{The Journal of Symbolic Logic} 74, no. 2 (2009):
  361--401. arXiv:math/0504196. DOI: 10.2178/jsl/1243948319.
\end{itemize}

\end{document}
